\documentclass[12pt]{article}
\usepackage{geometry}
\usepackage{setspace}
\usepackage{hyperref}
\usepackage{xcolor}
\usepackage{graphicx}
\usepackage{verbatim}
\geometry{margin=1in}
\setstretch{1.2}

\title{A Wave Curvature Equation for Unification \\ Part Two: The Impact of Waves}
\author{C.R. Kunferman}
\date{}

\begin{document}

\maketitle

\section*{Abstract}

Utilizing the calculations derived in Part 1 of this series, experiments and simulations are performed using various tools such as Python, Excel, and here on the Wolfram Alpha Notebook Cloud. The data produced presents impactful visual confirmation towards the hypothesis that there is a longitudinal pressure wave existent in the gravitational fields along with an extensive list of other implications. This paper explores these waves and their influence presented in simulation with some very impactful results.

\textbf{Note:} All charts and data presented here are simulated unless otherwise noted, the coding of which can be found by expanding the cells.

\section*{Hypothesis}

A longitudinal pressure wave produced by the excitation of the gravitational potential by gravitational waves, travelling at the speed of light, propagating over oscillating electric and magnetic fields transfers momentum and energy to matter, and the space time curvature.

\section*{Methods}

Under the assumption that the potential gravitational pressure wave herein "PGW" exists but has not been measured and under the observation of the similarity of wave calculations, simple experiments done with sound can be performed on different scales to represent the PGW.

It is well known and has been shown that a sound wave can transfer momentum and energy through various objects, mostly upon reaching a resonant frequency with the matter. These resonant waves can push down a wall, or shatter glass.

The first simulation was a combination of both, with very little complexity and outside influence, and shows an overly simplified example of how this is known to happen.

\subsection*{\textcolor{red}{In a real world setting the following materials would be needed:}}

\begin{itemize}
    \item A wave generator/loudspeaker with appropriate sound spectrum output within range of the object or material's resonant frequency, with controllable frequency and decibel controls, and an amplifier.
    \item An object made of glass in which the resonant frequency is known.
    \item Appropriate safety equipment, clothing, and protective wear for eyes, face, and body. Or a sealed testing room with bulletproof glass.
    \item A spectrometer.
    \item A decibel meter.
\end{itemize}

\subsection*{Experimental Setup:}

\begin{itemize}
    \item Place the glass object, wave generator, spectrometer, and decibel meter into a safe area, double checking that the surrounding areas are also safe in all directions, with a significant buffer for any unpredictable shards of glass that may exceed the testing area. Protective measures may need to be taken with cords and equipment as well.
    \item If it is contained within a safe room, take all controls outside of the area so you may operate them safely.
    \item Before connecting equipment, position the wave generator towards the glass object. In theory the object will either be knocked over when resonance is reached, or explode in all directions if it is exceeded, so the direction of the loudspeaker towards the object should be pointed away from the control area.
    \item Once the wave generator is positioned, proceed to cabling the spectrometer and decibel meter.
    \item Test both the spectrometer and decibel meter by yelling or clapping to make sure it is receiving input and providing readings.
    \item Remove the glass object, placing a marker as to where it belongs in the experiment.
    \item Connect the controls to the amplifier and set all volumes to negative values or 0 depending on the amplifier, making sure that it is not producing any output.
    \item Connect the wave generator to the amplifier.
    \item Send a test frequency through the amplifier and gradually increase volume to double check that both the spectrometer and decibel meter are reading the correct measurements with accuracy.
    \item Calibrate spectrometer and decibel meter as needed.
    \item If everything looks correct and you are ready to proceed, clear the staging area of all persons.
    \item Return amplifier to negative or 0 db, do not power down.
    \item Designate 1 researcher to enter the experiment area and replace the marker with the glass object. For added safety, an automated or remote control robotic may assume this role.
    \item Clear experiment area of all persons and any equipment not being used.
    \item Proceed to testing.
\end{itemize}

\section*{Experimental Testing}

\begin{itemize}
    \item Alert facility and surrounding areas that testing will be commencing.
    \item Test in intervals of decibels. Start with 1 db.
    \item Incrementally raise the frequency of the wave in steps from lowest to highest*, allowing enough time for the wave to propagate to the object using the calculation for the speed of sound, and a momentary buffer for the wave to travel through the object if resonance occurs.
    \item \textbf{\textcolor{red}{Warning:}} Do not exceed the resonant frequency of the safety equipment, or other equipment present within the testing area.
    \item If nothing occurs, raise the decibel by one, taking note of decibels tested. Repeat step 3.
\end{itemize}

\section*{Experimental Results}

\begin{itemize}
    \item When the wave reaches a particular decibel and frequency that is resonant with the object, it will exert energy onto it, effectively moving it in the direction of propagation.
    \item If the frequency exceeds the frequency, the vibration caused inside the object will cause it to expand in all directions, creating an explosion.
    \item With the information provided in both instances 1 and 2 from the spectrometer, you will have the resonant frequency and the threshold of exceeding the resonance.
\end{itemize}

\subsection*{Real World Application}

This experiment can be useful in a number of real world applications in various fields.

Examples include:

\begin{itemize}
    \item Medical: For breaking up kidney stones using ultrasound.
    \item Construction: For demolition, or for moving small objects.
\end{itemize}

\section*{Scaling Experiments with Calculations}

Translating the previous experiment to test and confirm the PGW requires a rethinking and new approach to the experiment, but contains essentially the same elements with different equipment. The overall goal being to generate pressure waves, and measure them.

\subsection*{Real World Experimental Setup for PGW}

\begin{itemize}
    \item A vacuum chamber capable of sustaining a low-pressure environment.
    \item A source of high-pressure gas, such as compressed air or a gas cylinder.
    \item A pressure wave generator (e.g., a piston or loudspeaker) capable of producing controlled pressure variations.
    \item Highly sensitive and precise instruments for measuring spacetime curvature, such as gravitational wave detectors or interferometers.
\end{itemize}

While we don’t have a generator currently for generating pressure waves traveling at the speed of light, it is necessary and also considered much safer to test this in simulation.

For purposes of presentation, the system state in the first experiment is set to more or less a perfect hypothetical state in which the resonance is equal to the pressure waves, and the pressure waves have a direct impact on the state metric. It can be thought of as a perfectly flat piece of fabric that exhibits the curvature and displacement of the pressure waves with little to no resistance given by its composition. This would be equivalent to thinly layered gas clouds or particles in space. For future refinement a differential equation, and calculations that account for the composition of the matter being influenced would be used.

Starting from this simple and easy to calculate primordial state provides a good starting point for observation that purifies the data as there will be no complexity added by outside influence, and will additionally keep it free of errors and contamination in testing.

To begin the simulations, a simple sine wave is generated that is resonant to the state metric. The simulation will calculate and plot its influence of the pressure wave on the state metric disk over time.

Having shown how a longitudinal pressure wave can be produced, and propagate, we can begin simulating a simplified version of the PGW’s influence, simplified as a sine wave:

\[
h(t,x) = \sin(2\pi t)
\]

\section*{Complexity}

While these simulated experiments are a simple example of how the PGW could affect both matter and the space time curvature, they lack real world data on a cosmological scale. However, these repeatable simulations have easily adjustable parameters and room for complexity to be added to test a wide gamut of frequencies. This provides a good foundation for the simulation and testing under different parameters.

\section*{Correlating the Solar System}

By applying real world data to the PGW, we can simulate, test, and refine our approach to see if there is any correlation between the data produced and observable effects taking place outside the simulation environments.

\textbf{Observations could include:}
\begin{itemize}
    \item Weather
    \item Surface Features
    \item Atmospheric Turbulence
    \item Tectonic Plates
    \item Planetary Accretion
    \item Orbiting Bodies/Moons Orbit
    \item Planetary Rings/Ring Position
    \item Planet Orbit
    \item Planetary Positioning and Angle of Rotation
    \item Anything seen as a result of gravity that gravitational waves cannot account for.
\end{itemize}

It is assumed that in a similar way shown in the experiments with sound waves that the PGW would have similar observable effects when reaching resonant frequencies on a much larger scale. To determine the resonant frequency of the PGW, the author explored different candidates and measurements to be used in order to determine the appropriate frequencies that this longitudinal pressure wave would be carrying in a cosmological environment.

The best candidate to begin this investigation would be the most readily observable effect or consequence of gravity, the \textbf{rate of gravity’s acceleration on the surface of a planet}. This would imply that the sine wave frequency would be equivalent to the rate of acceleration, and could be reverse calculated by making \(n\) equivalent. For clarity we can then replace the variable \(n\) with \(g\).

Expressed as:

\[
h(t,x) = \sin(2 \pi g t)
\]

where:

\(g\) is the frequency equivalent to the gravitational rate of acceleration on a planet’s surface.

We can begin our simulations by applying a range to \(t\), and calculating the rate of acceleration based on current measurements from reputable sources such as NASA.

Taking measurements and calculating the surface gravity of each planet in our solar system, the PGW for each were simulated over a period of 360\(t\), or range 0–360, to see if there were any noticeable and observable characteristics to strengthen this theory.

\section*{Planetary PGW Point Plots}

% Include planetary PGW point plots here as images, e.g.:
% \includegraphics[width=0.7\textwidth]{mercury_plot.png}
% \includegraphics[width=0.7\textwidth]{venus_plot.png}
% \includegraphics[width=0.7\textwidth]{uranus_plot.png}
% \includegraphics[width=0.7\textwidth]{pluto_plot.png}

\section*{Super Imposed PGW Point Plots on Surface Maps}

The data and charts show that the frequencies of each planet produce unique oscillatory patterns when viewed as scatter or point plots. Of particular note is Neptune, which oscillates in a very gradual and slow pattern, compared to a similar frequency on Earth which seems to oscillate quite frequently.

These charts were then combined on top of surface textures and relief maps to see if the data matched any of the observable real world features.

% Include surface map images here, e.g.:
% \includegraphics[width=0.7\textwidth]{mercury_surface.png}
% \includegraphics[width=0.7\textwidth]{venus_surface.png}

\section*{Experimental Result Analysis}

The analysis of the different plot styles reveals remarkable and unprecedented correlations with many of the planetary bodies in our solar system. These findings suggest a deep connection between oscillatory patterns and planetary features.

\subsection*{Surface Map Observations}

The surface maps provide striking visual evidence of these correlations. Unique oscillations, tied to each planet's frequencies, appear to shape observable planetary features across various scales—above ground, on the surface, and possibly even below.

\subsection*{Point Plot Insights}

Point plots further underscore these correlations, especially with planetary weather systems:

\begin{itemize}
    \item \textbf{Earth:} The points trace cloud paths and formation zones with remarkable precision.
    \item \textbf{Jupiter:} Data points align perfectly with the cloud bands and flank the iconic Great Red Spot, highlighting the storm’s influence.
    \item \textbf{Saturn:} The cloud bands are precisely delineated, with each data point occurring at the central axis of a band.
    \item \textbf{Neptune:} Data points cluster around a surface storm, seemingly intensifying its pressure, while the wave patterns drive weather systems at extreme speeds.
    \item \textbf{Uranus:} The patterns resemble Mars’ argyle-like formations but on a larger scale.
\end{itemize}

These alignments suggest that resonant gravity waves play a fundamental role in shaping weather and atmospheric structures.

\subsection*{Gravitational Acceleration as a Resonant Frequency}

The cornerstone of this investigation is the hypothesis that the sine wave frequency of the PGW can be equated to the rate of gravitational acceleration (\(g\)) on a planet’s surface. This idea is grounded in the premise that the most readily observable consequence of gravity is its acceleration rate, making it an ideal candidate for deriving the wave’s resonant frequency. By expressing the PGW as:

\[
h(t, x) = \sin(2 \pi g t),
\]

where \(g\) is the gravitational acceleration, the equation becomes directly tied to measurable planetary properties. This approach allows for simulations that use well-established datasets, such as those provided by NASA, to explore the effects of the PGW across different celestial environments.

\subsection*{Implications of the PGW Framework}

This framework has significant implications for understanding planetary systems:

\begin{enumerate}
    \item \textbf{Surface Oscillations and Patterns:} The alignment of oscillatory patterns with planetary surface features, such as Mars’ argyle-like formations and the separation of soil compositions, suggests that PGWs may influence the distribution and structure of surface materials. This raises the possibility that resonant frequencies uniquely shape each planet’s observable surface characteristics.
    \item \textbf{Atmospheric Dynamics:} The PGW framework provides a plausible explanation for the alignment between oscillatory patterns and atmospheric phenomena. On planets such as Jupiter and Saturn, the alignment of PGW data points with cloud bands and storm systems suggests that these waves may guide atmospheric flows and contribute to the formation of stable weather patterns.
    \item \textbf{Planetary Rings and Orbital Structures:} The appearance of gaps in contour plots, corresponding to ring systems on Saturn, Uranus, and Neptune, provides further evidence of the PGW’s influence. These observations suggest that resonant frequencies may play a role in organizing orbital debris and maintaining ring structures.
    \item \textbf{Universal Applicability:} While the PGW framework is rooted in gravitational acceleration as its foundational parameter, its universality extends to any celestial body where \(g\) can be measured. This makes the framework adaptable to exoplanets, moons, and even asteroid systems, offering a consistent method for exploring gravitational wave effects across the cosmos.
\end{enumerate}

\section*{Broader Scientific Context}

The introduction of the PGW as a driving force challenges traditional interpretations of planetary phenomena. Historically, features such as atmospheric bands, surface patterns, or ring structures have been studied as isolated phenomena with planet-specific explanations. The PGW framework unifies these observations under a single equation, providing a cohesive explanation that ties them to a fundamental cosmological force.

Additionally, the PGW framework aligns with experimental evidence from sound wave studies, where resonant frequencies produce observable patterns in physical media. Extending this concept to the cosmological scale represents a bold step in our understanding of gravitational forces and their consequences.

\section*{Conclusion}

The journey of discovery continues to surprise, and the charts and data produced from this simple equation are showing nearly definitive evidence that the PGW is a fundamental aspect of our universe, but it doesn’t end there.

\section*{Future Work}

In part 3 of this series, the author will showcase his work simulating the addition of Magnetic Fields to the PGW and will demonstrate the interplay of these forces on a cosmic scale, much larger than anything we have the ability to measure with our current technology.

The code above produces the following results:

\begin{verbatim}
Refined Earth-like Planets:
pI_name    MassKg       RadiusM     SurfaceGravity
PGWFrequency pI_orbsmax st_teff
-------------------------------------------------
['Kepler-197 b'] 6.2109e+24 6.4984e+06 9.8162
-0.91474 0.06 6004
['Kepler-192 d'] 6.2109e+24 6.4984e+06 9.8162
-0.91474 0.0686 5487
\end{verbatim}

\section*{Supplemental}

Further study was conducted of the PGW and while simulating earth's PGW in \texttt{Matlab}, the author found what appears to be a double helix within a 3D Swarm Chart. The implications of this are immense and warrant further study.

Additionally, this observation produced a unique way to search for life in the universe, under the assumption that a planet with a similar PGW and other habitable traits would have the same strand-like structure that may have guided the evolution of life on our planet.

By searching exoplanet data, the list of candidates has been narrowed down to a list of 6. The Matlab code to conduct this survey is as follows:

\begin{verbatim}
Constants
=6.67430e-11; % Gravitational constant in m^3 kg^-1 s^-2
rthMass = 5.972e24; % Earth's mass in kg
rthRadius = 6371e3; % Earth's radius in meters
\end{verbatim}

Subsequent or additional studies, derived works, and peer-reviewed edits are most welcome! Simply copy your notebooks to the author so he may assess them, address any questions or corrections, and ultimately add them with credit to this amazing exploration.

The author looks forward to collaboration within the scientific community, for the greater understanding of humanity.

\end{document}
